\phantomsection
\section*{Заключение}
\addcontentsline{toc}{section}{Заключение}

В результате выполнения курсового проекта разработано веб-приложение для
центра изучения иностранных языков, обеспечивающее запись студентов на
занятия, управление расписанием, публикацию учебных материалов и
отслеживание успеваемости в рамках единой системы с ролевым разграничением
доступа.

В ходе работы решены следующие \textbf{задачи}:

\begin{enumerate}
    \item Проведён анализ существующих образовательных платформ (Moodle,
    Google Classroom, iSpring Learn, Bitrix24) и обоснован выбор
    собственного решения на стеке Flask + PostgreSQL.
    \item Спроектирована реляционная модель данных, включающая десять
    сущностей и отражающая полный учебный процесс языкового центра.
    \item Реализована система аутентификации с ролевой моделью (студент,
    преподаватель, администратор) на основе Flask-Login и декораторов
    \texttt{@teacher\_required} / \texttt{@admin\_required}.
    \item Разработан каталог курсов с полнотекстовым поиском и фильтрацией
    по языку и уровню владения.
    \item Реализовано недельное расписание занятий с навигацией по неделям,
    фильтрами и записью с контролем вместимости.
    \item Добавлены учебные материалы с поддержкой файловых вложений
    (UUID-именование, скачивание через \texttt{send\_file}).
    \item Реализована система оценивания: преподаватель выставляет баллы
    и отметки посещаемости, студент просматривает агрегированный прогресс;
    предусмотрен экспорт журнала в формат CSV.
    \item Реализована система уведомлений с JS polling и toast-всплывашками.
    \item Настроено контейнерное развёртывание через Docker Compose
    с автоматическим применением Alembic-миграций.
\end{enumerate}

На основе полученных результатов можно сформулировать следующие
\textbf{выводы}:

\begin{enumerate}
    \item Использование Flask~3 с паттерном Application Factory обеспечивает
    чёткое разделение ответственности между модулями: каждый блупринт
    отвечает за свою область, что упрощает сопровождение и расширение
    приложения.
    \item PostgreSQL в сочетании с Flask-Migrate
    позволяет безопасно эволюционировать схему базы данных без потери
    данных, а SQLAlchemy исключает SQL-инъекции
    на уровне ORM.
    \item Контейнеризация с Docker гарантирует
    воспроизводимость среды выполнения и упрощает развёртывание
    как в локальной, так и в облачной инфраструктуре.
\end{enumerate}

\textbf{Перспективы дальнейшего развития.} В качестве возможных направлений
совершенствования приложения можно выделить: добавление мобильного
интерфейса (PWA или нативное приложение), интеграцию с внешними
календарными сервисами (Google Calendar, iCal), реализацию модуля
онлайн-оплаты занятий, а также внедрение push-уведомлений для замены
текущего JS polling на событийную модель через WebSocket.

При разработке пояснительной записки к курсовому проекту строго
соблюдались требования~\cite{gost_report}, определяющие структуру
и правила оформления научно-технических работ.

Исходный код проекта размещён в открытом репозитории:
\begin{center}
    \url{https://github.com/Miraz11287/foreign-language-center-webapp}.
\end{center}

Развёрнутая версия приложения доступна по адресу:
\begin{center}
    \url{https://langcenter.digitaldrugs.tech/}.
\end{center}
